% Step 3 research note: Ishida--LeBrun boundary and near saturation
% This note supersedes the over-broad conclusions and the false Weyl-only derivation in Step 2.

\section{Correction and live critical family}

The published mixed scalar--Weyl estimate is retained:
\[
 \frac1{4\pi^2}\int_M\left(\frac{s^2}{24}+2|W_+|^2\right)d\mu
 \geq \frac23(a^+)^2.
\]
The previous attempted deduction
\(\int|W_+|^2\ge(4\pi^2/3)(a^+)^2\) is withdrawn.  For a self-dual form \(\varphi\), the relevant Weitzenb\"ock estimate contains
\[
 \int(-W_+)(\varphi,\varphi)
 \geq -\frac16\int s|\varphi|^2-\frac12\int|\nabla\varphi|^2.
\]
The last term has the wrong sign to discard.  Moreover, \(F_A^+\) need not be harmonic and the pointwise spinor bound used in the 1998 argument is obtained only after passing to a negative constant-scalar-curvature Yamabe representative.  The genuine stronger primitive estimate is
\[
 \int_M(s-\sqrt6|W_+|)^2d\mu\ge72\pi^2(a^+)^2,
\]
and the mixed inequality follows from the exact completion of squares
\[
 27\int\left(\frac{s^2}{24}+2|W_+|^2\right)
 -\int(s-\sqrt6|W_+|)^2
 =\frac18\int(s+8\sqrt6|W_+|)^2.
\]

The critical counterexample candidates have therefore not disappeared.  Let
\[
 M=\left(\#_{j=1}^mX_j\right)\#N,\qquad m=2\ \hbox{or}\ 3,
\]
where the \(X_j\) are the symplectic summands used by Ishida--LeBrun, put
\[
 A=\sum_{j=1}^m c_1^2(X_j),\qquad
 D=4m-(2\chi+3\tau)(N).
\]
Then
\[
 T:=(2\chi+3\tau)(M)=A-D.
\]
Ishida--LeBrun exclude every Einstein metric when \(D\ge A/3\).  At the exact boundary
\[
 D=A/3,\qquad T=2A/3,
\]
their proof uses strictness for every individual metric.  It supplies no positive infimum over all metrics.  This is the sharp family on which to study whether
\(\inf_g\|g^{-1}(\operatorname{Ric}_g+g)\|_{L^\infty}\) can vanish.

\section{Exact necessary structure of a pinched boundary sequence}

Assume on one such boundary manifold (or on a varying sequence of such manifolds) that
\[
 \delta_i:=\|g_i^{-1}(\operatorname{Ric}_{g_i}+g_i)\|_{L^\infty}\longrightarrow0.
\]
Set
\[
 S_i=\frac1{4\pi^2}\int\frac{s_i^2}{24},\quad
 W_i=\frac1{4\pi^2}\int2|W_{+,i}|^2,\quad
 R_i=\frac1{4\pi^2}\int\frac{|\operatorname{Ric}_i^\circ|^2}{2},
\quad I_i=S_i+W_i.
\]
Chern--Gauss--Bonnet gives
\[
 I_i-R_i=\frac{2A}{3}.                                      \tag{1}
\]
If \(e_1,\dots,e_4\in[-\delta_i,\delta_i]\) are the eigenvalues of
\(g_i^{-1}(\operatorname{Ric}_{g_i}+g_i)\), then
\[
 |s_i+4|\le4\delta_i,
 \qquad |\operatorname{Ric}_i^\circ|^2
 =\sum(e_\alpha-\bar e)^2\le4\delta_i^2.
\]
Consequently
\[
 R_i\le q_iS_i,\qquad q_i=\frac{3\delta_i^2}{(1-\delta_i)^2}.
\]
The mixed monopole estimate gives \(I_i\ge2A/3\), and (1) implies
\[
 0\le R_i\le \frac{(2A/3)q_i}{1-q_i},\qquad
 I_i\to\frac{2A}{3}.                                      \tag{2}
\]
Here and below convergence is normalized by \(A\) if the topology varies.

Use all Bauer--Furuta sign choices
\[
 a_{\varepsilon,\eta}=\sum_j\varepsilon_jc_j+\sum_\ell\eta_\ell E_\ell.
\]
For every sign vector,
\[
 (a_{\varepsilon,\eta}^+)^2\le\frac32I_i=A+\frac32R_i.
\]
Averaging over signs cancels cross terms and yields
\[
 \operatorname{Av}(a_{\varepsilon,\eta}^+)^2
 =A+\sum_j\|c_{j,i}^-\|_2^2+\sum_\ell\|E_{\ell,i}^+\|_2^2.
\]
Thus
\[
 \sum_j\|c_{j,i}^-\|_2^2+\sum_\ell\|E_{\ell,i}^+\|_2^2
 \le\frac32R_i=O(A\delta_i^2).                            \tag{3}
\]
The period plane therefore asymptotically contains the canonical positive directions of all symplectic summands and becomes orthogonal to the negative directions of \(N\).  This is a substantive restriction on any degeneration, not merely an integral curvature statement.

Let
\[
 K_i=\frac1{72\pi^2}\int(s_i-\sqrt6|W_{+,i}|)^2.
\]
The scalar--Weyl bound gives \(K_i\ge(a^+)^2\), while the completion square gives
\(I_i\ge(2/3)K_i\).  Equations (2)--(3) force, for every fixed sign vertex,
\[
 K_i\to A,
 \qquad \frac1{864\pi^2}\int(s_i+8\sqrt6|W_{+,i}|)^2\to0. \tag{4}
\]
Since \(s_i\to-4\) uniformly, solving the two linear energy relations gives the unique profile
\[
 \boxed{R_i\to0,\quad S_i\to\frac{16A}{27},\quad
 W_i\to\frac{2A}{27},\quad K_i\to A}.                    \tag{5}
\]
In unnormalized notation,
\[
 \boxed{\operatorname{Vol}(M,g_i)\to\frac{32\pi^2A}{9}}, \tag{6}
\]
\[
 \boxed{\frac1{\pi^2A}\int|W_+|^2\to\frac4{27},\qquad
 \frac1{\pi^2A}\int|\operatorname{Ric}^\circ|^2\to0}.   \tag{7}
\]
In particular, the volume cannot escape on the exact boundary; it is automatically calibrated at order \(A\).  The reverse-orientation formula also fixes the total \(W_-\) energy, but does not prevent its concentration.

\section{Why known integral minimizers do not construct the sequence}

Ishida--LeBrun's exact \(L^2\)-Ricci infimum is attained only as a degenerating gluing sequence.  Its pieces are negative K\"ahler--Einstein orbifolds, scalar-flat anti-self-dual Green-function connectors, Ricci-flat gravitational instantons, and interpolation annuli.  For an exact scalar-flat region,
\[
 |s+4|=4\le4\delta,
\]
so \(\delta\ge1\), independently of how little volume that region has.  Hence the standard integral-minimizing neck cannot be a uniform almost-\(\operatorname{Ric}=-g\) neck.

This is also visible numerically.  Under the extra hypotheses of the 2003 connected-sum theorem, at \(D=A/3\),
\[
 \mathcal I_r=8\pi^2(A+D)=\frac{32\pi^2A}{3},
\]
whereas a hypothetical sequence satisfying (5)--(7) would have
\[
 \int|\operatorname{Ric}|^2\to\frac{128\pi^2A}{9},
\]
which is strictly larger.  It is therefore a different variational regime; exact knowledge of \(\mathcal I_r\) neither constructs nor excludes it.

At a vanishing geometric scale \(\ell_i\), rescale by \(\widehat g_i=\ell_i^{-2}g_i\).  Then
\[
 \operatorname{Ric}_{\widehat g_i}
 =-\ell_i^2\widehat g_i+O(\delta_i\ell_i^2)\widehat g_i.
\]
Every smooth noncollapsed blow-up limit is Ricci-flat.  Thus any bubble/cap in a pointwise pinched sequence must be Ricci-flat, whereas Burns and scalar-flat ASD connectors are not enough.  A complete noncollapsed separating connector with at least two ALE ends would be Ricci-flat; the splitting theorem then gives a line factor and linear growth, incompatible with four-dimensional ALE growth.  This excludes that specific multi-ended noncollapsed neck channel, not anisotropic collapse or cusp-like escape.

\section{The near-equality mechanism and its unresolved upgrade}

The exact equality proof has two stages.  Proposition 3.8 first turns equality of
\(\int(s-\sqrt6|W_+|)^2\) into a closed self-dual form of constant nonzero norm, lying in the lowest \(W_+\)-eigenline, with the other two eigenvalues equal and \(s+s^*\) a negative constant.  Convex-hull equality then aligns all relevant monopole vertices.  Finally the Blair-connection identities and Chern--Weil formula force \(s^*=s\), hence \(\nabla\omega=0\), and \(\operatorname{Ric}^\circ=0\), yielding a negative K\"ahler--Einstein metric.

Equations (2)--(4) show that every inequality in this chain is asymptotically saturated.  What is not in the source is a quantitative theorem upgrading this total deficit to a globally nonvanishing almost-K\"ahler form.  The missing statement should track the two H\"older steps, the rescaled Seiberg--Witten Weitzenb\"ock identity, the harmonic projection defect, and the conformal minimization.  Under a uniform Sobolev/harmonic-radius hypothesis it should imply, after gauge,
\[
 \|F_{A_i}^+-2\pi h_i^{a}\|_2\to0,
 \quad \|d\psi_i\|_2\to0,
 \quad ||\psi_i|-c|\to0,
\]
and eventually a positive pointwise lower bound for the harmonic representative.  That would produce a symplectic structure with the forbidden Chern class and rule out the boundary sequence.

Without such compactness, the natural alternative is phase separation.  Writing \(u_{j,i}\) for the harmonic representative of \(c_j^+\), all sign combinations
\(\sum_j\varepsilon_ju_{j,i}\) nearly saturate the same constant-length equality.  If quantitative equality holds locally, Walsh polarization forces
\[
 \sum_j|u_{j,i}|^2\approx c^2,
 \qquad \langle u_{p,i},u_{q,i}\rangle\approx0\quad(p\ne q).
\]
Because self-dual forms of a fixed orientation at one point form a three-dimensional Euclidean space and any two nonzero harmonic self-dual forms defining the same lowest eigenline are collinear, simultaneous near equality of distinct sign vertices cannot persist homogeneously.  The plausible escape is spatial localization: different canonical classes dominate on different macroscopic pieces, their zero/transition sets lying in collapsing or curvature-concentrating regions.  The connected-sum topology and the uncontrolled \(W_-\) energy would be carried precisely by those regions.

This yields the central dichotomy worth proving:
\begin{quote}
For an Ishida--LeBrun boundary sequence satisfying \(\delta_i\to0\), either the near-equality forms admit a uniform nonzero limit and create a forbidden symplectic/K\"ahler structure, or their mass phase-separates across a thin topology core.  In the latter case that core must be anisotropically collapsed or cusp-like; it cannot be the standard scalar-flat ASD connector and any noncollapsed bubble is Ricci-flat.
\end{quote}

The two decisive mathematical problems are therefore: (i) a quantitative near-equality theorem for the genuine scalar--Weyl estimate, with explicit dependence on Sobolev/harmonic radius; and (ii) a collapse/compactness theorem for \(C^0\)-almost negative Einstein four-metrics with the calibrated energy profile (3)--(7), describing whether phase-separated monopole forms can pass through a thin topology core.

\section{Relation to the complete Yau problem}

This boundary analysis addresses only the normalized negative branch in dimension four.  Even if it produced a counterexample sequence, a full refutation of Yau's interval statement would still require failure of the positive branch in the same dimension.  Conversely, proving a positive gap on these boundary families would be a major rigidity theorem but would not alone prove the full statement in all dimensions.  Its value is that it isolates the precise interface where a genuine all-Einstein obstruction and arbitrarily tight Ricci pinching might coexist, and turns the previous qualitative strictness into a sharply calibrated compactness-versus-degeneration problem.
